% Generated draft: requires mathematical and source review.
\documentclass[11pt]{article}
\usepackage[margin=25mm]{geometry}
\usepackage{fontspec}
\setmainfont{Latin Modern Roman}
\usepackage{amsmath,amssymb,mathtools,unicode-math}
\setmathfont{Latin Modern Math}
\usepackage{enumitem,needspace,xurl,hyperref}
\setlength{\parindent}{0pt}
\setlength{\parskip}{5pt}
\setlength{\emergencystretch}{3em}
\newcommand{\N}{\mathbb{N}}
\newcommand{\Z}{\mathbb{Z}}
\newcommand{\Q}{\mathbb{Q}}
\newcommand{\R}{\mathbb{R}}
\newcommand{\C}{\mathbb{C}}
\newcommand{\F}{\mathbb{F}}
\newcommand{\A}{\mathbb{A}}
\newcommand{\PP}{\mathbb{P}}
\newcommand{\E}{\mathbb{E}}
\newcommand{\eps}{\varepsilon}
\newcommand{\dd}{\,\mathrm{d}}
\begin{document}
\section*{UnsolvedMath ID: 700444}
\textbf{Source problem code:} catalogue entry
\subsection{Definitions and mathematical statement}
Let \(\Omega\) be a nonempty set, finite or infinite, and let \(G\leq \operatorname{Sym}(\Omega)\) be a permutation group. The action is transitive if for all \(\alpha,\beta\in\Omega\) there exists \(g\in G\) with \(g(\alpha)=\beta\). A transitive action is primitive if its only \(G\)-invariant equivalence relations on \(\Omega\) are equality and the universal relation. For pairwise distinct \(\alpha,\beta,\gamma\in\Omega\), define the pointwise stabilizer by
\[
G_{\alpha,\beta,\gamma}=\{g\in G:g(\alpha)=\alpha,\ g(\beta)=\beta,\ g(\gamma)=\gamma\}.
\]
The problem asks to describe the structure of all primitive permutation groups \(G\leq\operatorname{Sym}(\Omega)\), in both finite and infinite degree, such that
\[
\forall\alpha,\beta,\gamma\in\Omega,\qquad
\bigl(\alpha\ne\beta\ \land\ \alpha\ne\gamma\ \land\ \beta\ne\gamma\bigr)
\Longrightarrow G_{\alpha,\beta,\gamma}=\{1_G\}.
\]
Here the stabilizer is interpreted as fixing each of the three points individually. The source gives no lower bound on \(|\Omega|\); the displayed condition is vacuous when \(|\Omega|<3\). The definition of primitivity above includes the singleton action, although conventions excluding degree one are also used.

The source additionally states that this problem is closely connected with describing groups having a Frobenius group as a maximal subgroup. This is preserved as a related question, not as an asserted equivalence. The source does not specify the precise connection or the convention for Frobenius groups in the infinite setting.
\subsection{Short English statement}
Describe the structure of finite and infinite primitive permutation groups in which the identity is the only element fixing three pairwise distinct points individually. The source also mentions a connection with groups having a Frobenius maximal subgroup, without specifying that connection.
\subsection{Sources}
\begin{itemize}
\item Local UnsolvedMath registry, permanent ID \(700444\), exact source code \texttt{catalogue entry}, title \emph{Kourovka Notebook Problem 7.50}. The supplied record labels the problem open; this status has not been independently verified.
\item A. N. Fomin, proposer of Problem 7.50, \emph{The Kourovka Notebook (Unsolved Problems in Group Theory)}, 7th issue (1980). The supplied background also identifies the 21st issue edition (2026), \texttt{arXiv:1401.0300v48}. The supplied collection description names E. I. Khukhro and V. D. Mazurov as editors.
\end{itemize}
\end{document}
